\scanpage{23}
\mspage{21}

\[
 \mathrm{Hom\,cat}_k(a,b)\overset{\unsure{\mathbb L|L}}{\otimes}_k\mathrm{Hom\,cat}_k(b,c)\mathbin{\unsure{-|\longrightarrow}}\mathrm{Hom\,cat}_k(a,c),
\]

d'où on déduira des accouplements plus faibles dans $\mathrm{Hot}_{\bullet}$ (au lieu de $\mathrm D_{\geq 0}(k)$)

\[
 \mathrm{Hom\,cat}_{\mathrm{Hot}_{\bullet}}(a,b)\times
 \mathrm{Hom\,cat}_{\mathrm{Hot}_{\bullet}}(b,c)\longrightarrow
 \mathrm{Hom\,cat}_{\mathrm{Hot}_{\bullet}}(a,c).
\]

$3^\circ)$ Revenant au cas $M$ ponctuée sans plus, Quillen définit, pour $a,b$ dans $M$

\[
 \unsure{\Pi|\pi}_{\unsure{1}}(a,b)=\Pi_1(a,b;0_a,0_b)
\]

\margin{comme bifoncteur sur $(M\Sigma^{-1})^{\unsure{\circ|0}}\times M\Sigma^{-1}$}

(avec les notations du $1^\circ)$). Si ce que j'avance sans preuves dans $1^\circ)$ est vrai, il doit s'ensuivre qu'on a un isom. can.

\[
 \Pi_1(a,b)\simeq\Pi_{\unsure{1|2}}\bigl(\mathrm{Ch}^{\Sigma}_{\infty}(a,b)\bigr)
\]

où $\mathrm{Ch}^{\Sigma}_{\infty}(a,b)$ est vu comme objet de $\mathrm{Hot}_{\bullet}$, donc son $\Pi_1$ est défini\unsure{.|:} On a \deleted{\unsure{ainsi}} \inserted{en tous cas} sans aucune référence à des cofibrations et des fibrations, un bifoncteur à valeurs dans Groupes, composé \deleted{des}

\diagram{\begin{tikzcd}[ampersand replacement=\&,column sep=large]
(M\Sigma^{-1})^{\circ}\times M\Sigma^{-1}
\arrow[r,"{(a,b)\mapsto\mathrm{Ch}^{\Sigma}_{\infty}(a,b)}"]\&
\mathrm{Hot}_{\bullet}\arrow[r,dash,"\Pi_1"]\&(\mathrm{Groupes})
\end{tikzcd}}

Si pour $a$ fixé, $b$ variable, \deleted{$\Pi_1(a,b)$}

\[
 b\mathbin{\unsure{-|\mapsto}}\Pi_1(a,b)
\]

est représentable, on note l'objet qui le représente par $\Sigma(a)$ \deleted{\unsure{$f$}} (\uline{suspension} -- ne pas confondre ce $\Sigma$ avec $\Sigma\subset\mathrm{Fl}(M)$, qu'il vaut mieux maintenant noter $W$, comme ``weak \unsure{equivalence|equivalences}''). Donc
% archive sheet number: 22
